% =============================================================================
%  Guide d'installation — Formation Python pour l'analyse de données
%
%  COMPILATION
%    Moteur   : pdfLaTeX  (fonctionne aussi avec XeLaTeX / LuaLaTeX)
%    VS Code  : extension « LaTeX Workshop », recette « pdflatex »
%    Passes   : 2 (la seconde pour les liens hyperref)
%
%  PAQUETS REQUIS (tous présents dans MiKTeX / TeX Live complet ;
%  MiKTeX les installera à la volée au premier passage)
%    babel, fontenc, helvet, geometry, xcolor, tcolorbox, fancyhdr,
%    enumitem, tabularx, pifont, fancyvrb, hyperref
%
%  PERSONNALISATION RAPIDE
%    - Couleurs      : bloc « Palette » ci-dessous
%    - Police        : remplacer helvet par lato, inter, roboto…
%    - Coordonnées   : chercher « CONTACT » en fin de document
% =============================================================================

\documentclass[11pt,a4paper]{article}

% --- Encodage et langue ------------------------------------------------------
\usepackage[T1]{fontenc}
\usepackage[utf8]{inputenc}      % inoffensif sur les moteurs modernes
\usepackage[french]{babel}

% --- Police sans empattement -------------------------------------------------
\usepackage{helvet}
\renewcommand{\familydefault}{\sfdefault}

% --- Mise en page ------------------------------------------------------------
\usepackage[a4paper,left=2.1cm,right=2.1cm,top=2.6cm,bottom=2.2cm,
            headheight=26pt,headsep=14pt,footskip=26pt]{geometry}
\usepackage{fancyhdr}
\usepackage{enumitem}
\usepackage{tabularx}
\usepackage{array}
\usepackage{pifont}
\usepackage{fancyvrb}
\usepackage{xcolor}
\usepackage[most]{tcolorbox}
\usepackage[hidelinks]{hyperref}
\hypersetup{
  pdftitle={Guide d'installation d'Anaconda — Formation Python pour l'analyse de données},
  pdfauthor={Fiacre Bandaogo},
  pdfsubject={Installation d'Anaconda et de JupyterLab, à faire avant la séance 1},
  pdfkeywords={Anaconda; JupyterLab; conda; Python; installation; formation}}

% =============================================================================
%  Palette
% =============================================================================
\definecolor{bleufonce}{HTML}{1B3A5C}
\definecolor{bleuclair}{HTML}{2E6DA4}
\definecolor{bleupale} {HTML}{9FC0DC}
\definecolor{grisfonce}{HTML}{333333}
\definecolor{grismoyen}{HTML}{6B7280}
\definecolor{grisclair}{HTML}{F4F6F8}
\definecolor{bordure}  {HTML}{DDE3E9}
\definecolor{orange}   {HTML}{C05621}
\definecolor{orangebg} {HTML}{FDF3E7}
\definecolor{vert}     {HTML}{2F6F4E}
\definecolor{vertbg}   {HTML}{EDF6F0}
\definecolor{rouge}    {HTML}{B02A2A}
\definecolor{rougebg}  {HTML}{FDECEC}
\definecolor{codebg}   {HTML}{EAF0F6}

\color{grisfonce}

% Pas d'alinéa : espacement vertical entre paragraphes
\setlength{\parindent}{0pt}
\setlength{\parskip}{6pt}

% =============================================================================
%  Listes
%  (babel-french remplace les puces par des tirets : on rétablit)
% =============================================================================
\setlist[itemize]{label=\textcolor{bleuclair}{\textbullet}, leftmargin=1.5em,
                  itemsep=4pt, topsep=5pt, parsep=0pt}
\setlist[enumerate]{label=\textbf{\arabic*}, leftmargin=1.9em,
                    itemsep=5pt, topsep=5pt, parsep=0pt}

% =============================================================================
%  En-tête et pied de page
% =============================================================================
\fancypagestyle{courant}{%
  \fancyhf{}
  \renewcommand{\headrulewidth}{0pt}
  \renewcommand{\footrulewidth}{0.4pt}
  \renewcommand{\footrule}{\hbox to\headwidth{\color{bordure}\leaders\hrule height 0.4pt\hfill}}
  \fancyhead[L]{%
    {\setlength{\fboxsep}{0pt}%
     \colorbox{bleufonce}{\makebox[\textwidth][l]{%
       \rule[-7pt]{0pt}{22pt}%
       \hspace{6pt}\footnotesize\bfseries\color{bleupale}%
       FORMATION PYTHON — GUIDE D'INSTALLATION\hfill
       À FAIRE AVANT LA SÉANCE 1\hspace{6pt}}}}}
  \fancyfoot[L]{\scriptsize\color{grismoyen}Formation Python pour l'analyse de données}
  \fancyfoot[R]{\scriptsize\color{grismoyen}Page \thepage}
}
\fancypagestyle{couverture}{%
  \fancyhf{}
  \renewcommand{\headrulewidth}{0pt}
  \renewcommand{\footrulewidth}{0.4pt}
  \renewcommand{\footrule}{\hbox to\headwidth{\color{bordure}\leaders\hrule height 0.4pt\hfill}}
  \fancyfoot[L]{\scriptsize\color{grismoyen}Formation Python pour l'analyse de données}
  \fancyfoot[R]{\scriptsize\color{grismoyen}Page \thepage}
}
\pagestyle{courant}

% =============================================================================
%  Titres
% =============================================================================
\newcommand{\titregrand}[1]{%
  \par\vspace{10pt}%
  {\Large\bfseries\color{bleufonce}#1}%
  \par\vspace{6pt}}

\newcommand{\titrepetit}[1]{%
  \par\vspace{10pt}%
  {\large\bfseries\color{bleuclair}#1}%
  \par\vspace{4pt}}

% --- Bandeau d'étape : \etape{numéro}{titre}{durée} --------------------------
\newcommand{\etape}[3]{%
  \par\vspace{12pt}%
  \noindent{\setlength{\fboxsep}{0pt}%
   \colorbox{bleufonce}{%
     \makebox[\textwidth][l]{%
       \rule[-8pt]{0pt}{26pt}%
       \hspace{7pt}%
       \parbox{2.0cm}{\scriptsize\bfseries\color{bleupale}ÉTAPE #1}%
       {\large\bfseries\color{white}#2}%
       {\footnotesize\color{bleupale}~ #3}\hfill}}}%
  \par\vspace{8pt}}

% =============================================================================
%  Encadrés
% =============================================================================
\newtcolorbox{encadre}[3]{%
  enhanced, breakable, sharp corners,
  colback=#2, colframe=#2, boxrule=0pt,
  borderline west={2.5pt}{0pt}{#1},
  left=9pt, right=9pt, top=7pt, bottom=7pt,
  before skip=8pt, after skip=8pt,
  fontupper=\small,
  title={}, before upper={\textcolor{#1}{\textbf{#3}}\par\vspace{3pt}}}

\newenvironment{info}[1]     {\begin{encadre}{bleuclair}{grisclair}{\ding{110}~~#1}}{\end{encadre}}
\newenvironment{attention}[1]{\begin{encadre}{orange}{orangebg}{\ding{110}~~#1}}{\end{encadre}}
\newenvironment{succes}[1]   {\begin{encadre}{vert}{vertbg}{\ding{51}~~#1}}{\end{encadre}}
\newenvironment{alerte}[1]   {\begin{encadre}{rouge}{rougebg}{\ding{110}~~#1}}{\end{encadre}}

% --- Encadré « problème » de la section dépannage ---------------------------
\newtcolorbox{probleme}[1]{%
  enhanced, breakable, sharp corners,
  colback=white, colframe=bordure, boxrule=0.6pt,
  left=9pt, right=9pt, top=7pt, bottom=7pt,
  before skip=6pt, after skip=6pt,
  fontupper=\small,
  before upper={\textcolor{bleufonce}{\textbf{#1}}\par\vspace{3pt}}}

% --- Bloc de code -----------------------------------------------------------
\newtcolorbox{codebloc}[1][codebg]{%
  enhanced, sharp corners,
  colback=#1, colframe=#1, boxrule=0pt,
  borderline west={2pt}{0pt}{bordure},
  left=10pt, right=8pt, top=5pt, bottom=5pt,
  before skip=6pt, after skip=6pt}

% --- Case à cocher ----------------------------------------------------------
\newcommand{\caseacocher}{%
  \textcolor{bleuclair}{\framebox[3.4mm]{\rule{0pt}{3.0mm}}}}

% --- CONTACT ----------------------------------------------------------------
%  UNE SEULE LIGNE À MODIFIER : remplacez l'adresse ci-dessous par la vôtre.
%  Elle est reprise automatiquement aux 6 endroits du guide où l'on demande
%  au participant d'écrire, et devient un lien mailto cliquable.
\newcommand{\courriel}{bandaogofiacre0@gmail.com}
\newcommand{\adressecontact}{%
  \href{mailto:\courriel}{\textbf{\color{bleuclair}\courriel}}}

% --- Raccourcis -------------------------------------------------------------
\newcommand{\fichier}[1]{\texttt{#1}}
\newcommand{\fleche}{~$\rightarrow$~}

% =============================================================================
\begin{document}
\thispagestyle{couverture}

% -----------------------------------------------------------------------------
%  Couverture
% -----------------------------------------------------------------------------
\noindent\begin{tcolorbox}[enhanced, sharp corners, colback=bleufonce,
  colframe=bleufonce, boxrule=0pt, left=14pt, right=14pt, top=14pt, bottom=14pt,
  before skip=0pt, after skip=16pt]
  {\footnotesize\bfseries\color{bleupale}%
   F\,O\,R\,M\,A\,T\,I\,O\,N\quad P\,Y\,T\,H\,O\,N\quad P\,O\,U\,R\quad
   L'A\,N\,A\,L\,Y\,S\,E\quad D\,E\quad D\,O\,N\,N\,É\,E\,S\par}
  \vspace{10pt}
  {\fontsize{25}{29}\selectfont\bfseries\color{white}Guide d'installation\par}
  \vspace{6pt}
  {\large\color{bleupale}Anaconda + JupyterLab — à réaliser avant la séance 1\par}
  \vspace{10pt}
  {\color{bleuclair}\hrule height 1pt}
  \vspace{8pt}
  {\small\color{bleupale}Durée estimée : 30 à 45 minutes\quad\textbullet\quad
   Aucune connaissance préalable requise\par}
\end{tcolorbox}

\noindent Ce guide vous accompagne pas à pas dans l'installation des outils
nécessaires à la formation.

\begin{attention}{Ne repoussez pas cette installation à la veille de la séance}
Environ un participant sur trois rencontre une difficulté, c'est normal et cela
se règle facilement, \textbf{à condition de s'y prendre à l'avance}. Faites
l'installation dès réception de ce guide. En cas de blocage,
écrivez à \adressecontact : on vous aidera à résoudre le problème avant
la première séance.
\end{attention}

\titrepetit{Les 5 étapes en un coup d'œil}

{\renewcommand{\arraystretch}{1.5}%
\noindent\begin{tabularx}{\textwidth}{@{}p{0.8cm}X@{}}
\hline
\textbf{1} & Télécharger Anaconda\hfill{\small\color{grismoyen}5 min + téléchargement}\\
\hline
\textbf{2} & Installer Anaconda\hfill{\small\color{grismoyen}15 min}\\
\hline
\textbf{3} & Enregistrer le dossier de la formation\hfill{\small\color{grismoyen}3 min}\\
\hline
\textbf{4} & Lancer la formation et vérifier\hfill{\small\color{grismoyen}20 min, presque sans rien faire}\\
\hline
\textbf{5} & Envoyer votre confirmation\hfill{\small\color{grismoyen}2 min}\\
\hline
\end{tabularx}}

\newpage

% -----------------------------------------------------------------------------
\etape{1}{Télécharger Anaconda}{}

Anaconda est une distribution qui installe d'un seul coup Python et tous les
outils dont nous aurons besoin pendant la formation.

\begin{enumerate}
\item Ouvrez le lien suivant :
      \textbf{\href{https://www.anaconda.com/download/success}%
      {\color{bleuclair}\underline{www.anaconda.com/download/success}}}
\item En haut de la page, sélectionnez l'onglet correspondant à votre
      ordinateur : \textbf{Windows}, \textbf{Mac} ou \textbf{Linux}.
\item Repérez le bloc intitulé \textbf{Anaconda Distribution}, et non le bloc
      \og Miniconda\fg{} situé en dessous.
\item Dans ce bloc, cliquez sur \textbf{Graphical Installer} pour démarrer le téléchargement.
\end{enumerate}

\begin{info}{Quelle version pour un Mac ?}
La page propose l'installateur correspondant aux Mac récents (puce Apple M1 à
M4). Pour connaître votre modèle : menu \textbf{Pomme}\fleche\textbf{À propos de
ce Mac}. Si vous avez un Mac plus ancien à processeur \textbf{Intel} et qu'aucune
version ne semble correspondre, signalez-le avant d'installer quoi que ce soit.
\end{info}

\begin{info}{Le fichier est volumineux}
Comptez environ 1 Go et 5 à 20 minutes de téléchargement selon votre connexion.
\end{info}

% -----------------------------------------------------------------------------
\etape{2}{Installer Anaconda}{}

\titrepetit{Sur Windows}

\begin{enumerate}
\item Ouvrez le fichier téléchargé.
\item Si Windows affiche un avertissement de sécurité, cliquez sur
      \textbf{Plus d'informations} puis \textbf{Exécuter quand même}.
\item Cliquez sur \textbf{Next}, puis \textbf{I Agree} pour accepter la licence.
\item À la question \og Install for\fg{}, laissez \textbf{Just Me (recommended)}.
\item Laissez le dossier d'installation proposé par défaut. Cliquez sur
      \textbf{Next}.
\item \textbf{Écran des options avancées, le seul écran qui demande votre
      attention :}

\vspace{2pt}
\noindent\begin{tcolorbox}[enhanced, sharp corners, colback=grisclair,
  colframe=bordure, boxrule=0.6pt, left=8pt, right=8pt, top=6pt, bottom=6pt,
  before skip=4pt, after skip=8pt]
{\renewcommand{\arraystretch}{1.35}%
\begin{tabularx}{\textwidth}{@{}p{0.6cm}X@{}}
\textcolor{rouge}{\ding{55}} & \small\og Add Anaconda3 to my PATH environment variable\fg{}%
  \hfill\textbf{\color{rouge}NE PAS COCHER}\\
\textcolor{vert}{\ding{51}}  & \small\og Register Anaconda3 as my default Python\fg{}%
  \hfill\textbf{\color{vert}LAISSER COCHÉ}\\
\end{tabularx}}
\end{tcolorbox}

\item  Cliquez sur \textbf{Install} et attendez que
l'installation se termine.

\vspace{4pt}
\item  À la fin, décochez les deux cases proposant de découvrir
Anaconda dans le navigateur, puis cliquez sur \textbf{Finish}.
\end{enumerate}

\newpage
\titrepetit{Sur macOS}

\begin{enumerate}
\item Ouvrez le fichier \fichier{.pkg} téléchargé.
\item Suivez l'assistant en cliquant sur \textbf{Continuer} à chaque écran, puis
      \textbf{Accepter} pour la licence.
\item À l'écran \og Type d'installation\fg{}, choisissez \textbf{Installer pour
      moi uniquement} si l'option est proposée.
\item Saisissez votre mot de passe de session si macOS le demande.
\item Attendez la fin de l'installation, puis cliquez sur \textbf{Fermer}.
\end{enumerate}

\begin{attention}{Poste professionnel : droits administrateur}
Si l'installation est refusée faute de droits administrateur, ne perdez pas de
temps à insister : écrivez immédiatement à \adressecontact. Il existe une solution de
contournement (voir la dernière page) et votre service informatique peut être
sollicité, ce qui prend souvent plusieurs jours.
\end{attention}

% -----------------------------------------------------------------------------
\etape{3}{Enregistrer le dossier de la formation}{}

Le dossier de la formation est en accès libre : vous le téléchargez vous-même,
sans rien demander à personne. Il contient votre dossier de travail et les
programmes qui configureront Python.

\begin{enumerate}
\item Ouvrez le lien suivant :
      \textbf{\href{https://github.com/Fiacrebandaogo/formation-python}%
      {\color{bleuclair}\underline{github.com/Fiacrebandaogo/formation-python}}}
\item Cliquez sur le bouton vert \textbf{Code}, au-dessus de la liste des
      fichiers, puis sur \textbf{Download ZIP} dans le menu qui s'ouvre.
\item \textbf{Décompressez} le fichier téléchargé : clic droit sur le fichier
      \fleche\textbf{Extraire tout} (Windows) ou double-clic (Mac).
\item Vous obtenez un dossier nommé \fichier{\textbf{formation-python-main}}.
      Placez-le où vous voulez sur votre ordinateur, mais retenez où.
\end{enumerate}

\begin{info}{Pourquoi \og{}-main\fg{} à la fin du nom ?}
C'est le site qui ajoute ce suffixe automatiquement. Ce n'est pas une erreur et
cela ne gêne en rien. Vous pouvez renommer le dossier en
\fichier{formation\_python} si vous préférez.
\end{info}

À l'intérieur se trouve un sous-dossier nommé
\fichier{\textbf{formation\_python}} : c'est celui dont vous aurez besoin à
l'étape suivante, et à chaque séance.

% -----------------------------------------------------------------------------
\etape{4}{Lancer la formation et vérifier}{}

Un seul fichier à retenir, et c'est celui que vous rouvrirez à chaque séance.
Ouvrez le sous-dossier \fichier{\textbf{formation\_python}}, à l'intérieur du
dossier téléchargé : le fichier à lancer s'y trouve. Au premier lancement, il
installe tous les outils avant d'ouvrir un éditeur de code appelé ``JupyterLab''.

\vspace{4pt}
\noindent\begin{tcolorbox}[enhanced, sharp corners, colback=grisclair,
  colframe=bordure, boxrule=0.6pt, left=8pt, right=8pt, top=6pt, bottom=6pt,
  before skip=4pt, after skip=8pt]
\begin{tabularx}{\textwidth}{@{}p{2.0cm}X@{}}
\textbf{Windows} & Double-cliquez sur \fichier{\textbf{lancer\_jupyterlab.bat}}\\[5pt]
\textbf{Mac}     & Double-cliquez sur \fichier{\textbf{lancer\_jupyterlab.command}}\\
\end{tabularx}
\end{tcolorbox}

Une fenêtre noire s'ouvre et du texte défile pendant 5 à 15 minutes.
C'est normal. \textbf{Ne fermez pas cette fenêtre et ne touchez à rien}.
JupyterLab s'ouvrira ensuite dans votre navigateur.

\begin{attention}{Cette fenêtre noire reste ouverte}
Elle fait tourner JupyterLab en arrière-plan : \textbf{ne la fermez pas} tant que
vous travaillez. Pour arrêter en fin de séance, il suffit de la fermer.
\end{attention}

\begin{info}{Un jour, ce fichier affichera \og MISE À JOUR\fg{}}
C'est normal : cela signifie que vous avez ajouté un outil (package) à la formation.
Laissez-le faire et JupyterLab s'ouvrira ensuite comme d'habitude.
\end{info}

\begin{info}{Sur Mac : \og le développeur ne peut pas être vérifié\fg{}}
Message normal : Faites un \textbf{clic droit} sur
\fichier{lancer\_jupyterlab.command}\fleche \textbf{Ouvrir}\fleche\textbf{Ouvrir}.
\end{info}

\begin{attention}{Si la fenêtre affiche \og INSTALLATION INTERROMPUE\fg{}}
Ce n'est ni grave ni définitif. Un fichier \fichier{installation.log} est
créé dans le dossier \fichier{formation\_python} :
\textbf{envoyez-le à \adressecontact}. Il contient la cause exacte et se règle en
général en quelques minutes.
\end{attention}

\titrepetit{Le test}
Après l'ouverture de JupyterLab, vous devez vérifier que tout fonctionne.
Pour cela, vous allez exécuter un petit programme de test.

\begin{enumerate}
\item Dans le panneau de gauche, double-cliquez sur le dossier
      \fichier{seances}.
\item Dans la zone principale, sous \textbf{Notebook}, cliquez sur
      \textbf{Python 3 (formation)}.
\item Un document s'ouvre avec une case vide : c'est une \textbf{cellule}.
      Cliquez dedans et recopiez ces deux lignes :

\begin{codebloc}
\begin{Verbatim}[fontsize=\small,commandchars=\\\{\}]
import pandas as pd
print('Tout fonctionne, version de pandas :', pd.__version__)
\end{Verbatim}
\end{codebloc}

\item Appuyez sur \textbf{Maj + Entrée} pour exécuter la cellule.
\end{enumerate}

Vous devez voir apparaître sous la cellule :

\begin{codebloc}[vertbg]
\begin{Verbatim}[fontsize=\small]
Tout fonctionne, version de pandas : 2.2.3
\end{Verbatim}
\end{codebloc}

Le numéro de version peut différer, c'est sans importance.

\begin{attention}{Un message rouge est apparu ?}
Fermez JupyterLab, puis double-cliquez sur
\fichier{\textbf{diagnostic.bat}} (ou \fichier{diagnostic.command} sur Mac). Ce
programme ne modifie rien : il crée un fichier \fichier{diagnostic.txt} à envoyer
à \adressecontact.
\end{attention}

% -----------------------------------------------------------------------------
\etape{5}{Envoyer votre confirmation}{}

Cette confirmation permet de repérer les participants en difficulté
\textbf{avant} la première séance, et de préparer une solution individuelle si
nécessaire.

\noindent Envoyez à \adressecontact un e-mail contenant :

\begin{itemize}
\item Une \textbf{capture d'écran} de votre notebook montrant
      \og Tout fonctionne\fg{}
\item La mention de toute difficulté rencontrée, même résolue
\end{itemize}

\vspace{-20pt}
\titrepetit{Et ensuite ?}

Vous n'avez rien à préparer d'autre : aucune lecture, aucun exercice préalable.
La séance 1 part de zéro et le premier code que vous écrirez sera écrit ensemble,
en salle. Apportez simplement votre ordinateur, chargé, avec le chargeur.


\newpage
% -----------------------------------------------------------------------------
\titregrand{Problèmes fréquents et solutions}
{\small\color{grismoyen}Ces situations sont courantes. Aucune n'est bloquante.\par}
\vspace{6pt}

\begin{probleme}{L'installation d'Anaconda est refusée : \og droits administrateur requis\fg{}}
Fréquent sur les ordinateurs professionnels. Écrivez à \adressecontact
immédiatement — le délai de traitement par un service informatique se compte en
jours. En attendant, utilisez la solution de secours ci-dessous.
\end{probleme}

\begin{probleme}{L'antivirus bloque l'installation d'Anaconda}
Suspendez temporairement l'antivirus le temps de l'installation, puis
réactivez-le. Si vous ne pouvez pas le faire, passez par votre service
informatique.
\end{probleme}

\begin{probleme}{Double-cliquer sur lancer\_jupyterlab.bat ne fait rien, ou la fenêtre se ferme aussitôt}
Le dossier n'a pas été décompressé : sous Windows, les programmes ne s'exécutent
pas depuis un fichier .zip. Faites un clic droit sur
\fichier{formation\_python.zip}\fleche\textbf{Extraire tout}, puis relancez
depuis le dossier obtenu.
\end{probleme}

\begin{probleme}{\og INSTALLATION INTERROMPUE\fg{} s'affiche}
Envoyez à \adressecontact le fichier \fichier{installation.log} créé dans le dossier
\fichier{formation\_python}. Les causes les plus fréquentes sont l'absence de
connexion, un réseau d'entreprise qui filtre les téléchargements, ou un manque
d'espace disque (prévoir 3 Go).
\end{probleme}

\begin{probleme}{\og Anaconda est introuvable sur cet ordinateur\fg{}}
Les étapes 1 et 2 n'ont pas abouti, ou Anaconda a été installé à un endroit
inhabituel. Reprenez ces deux étapes, puis relancez
\fichier{lancer\_jupyterlab.bat}.
\end{probleme}

\begin{probleme}{Sur Mac : \og le développeur ne peut pas être vérifié\fg{}}
Normal pour un fichier reçu par e-mail. Clic droit sur le fichier\fleche
\textbf{Ouvrir}\fleche\textbf{Ouvrir}. À ne faire qu'une fois par fichier.
\end{probleme}

\begin{probleme}{\og No module named ...\fg{} dans un notebook, alors que l'installation a réussi}
Un ancien noyau Jupyter présent sur votre machine a pris le dessus. Dans
JupyterLab : \textbf{Kernel}\fleche\textbf{Change Kernel}\fleche
\og \textbf{Python 3 (formation)}\fg{}, puis réexécutez la cellule.
\end{probleme}

\begin{probleme}{Autre problème, ou vous ne savez pas quoi faire}
Double-cliquez sur \fichier{\textbf{diagnostic.bat}}
(\fichier{diagnostic.command} sur Mac). Il ne modifie rien et produit un fichier
\fichier{diagnostic.txt} : envoyez-le à \adressecontact, il contient tout ce qu'il
faut savoir.
\end{probleme}

\vspace{6pt}
\begin{alerte}{Solution de secours : Google Colab}
Si l'installation s'avère impossible sur votre poste, \textbf{vous ne serez pas
pénalisé}. Google Colab permet de suivre l'intégralité de la formation depuis un
simple navigateur, sans rien installer : il suffit d'un compte Google et de se
rendre sur \textbf{\href{https://colab.research.google.com}%
{\color{rouge}\underline{colab.research.google.com}}}.

\vspace{4pt}
Écrivez à \adressecontact pour qu'on vous transmette les supports au format adapté.
En parallèle, nous continuerons à chercher une solution pour installer Anaconda
sur votre poste.
\end{alerte}

\end{document}
